%%%%%%%%%%%%%%%%%%%%%%%%%%%%%%%%%%%%%%%%%%%%%%%%%%%%%%%%%%%%%%%%%%%%%%%%%%%
% Chapter: From Sensor to Digital Value (capstone)
% Falstad examples: examples/capstone (falstad/circuits/capstone.py)
% All numbers in this chapter are checked by capstone.py.
%%%%%%%%%%%%%%%%%%%%%%%%%%%%%%%%%%%%%%%%%%%%%%%%%%%%%%%%%%%%%%%%%%%%%%%%%%%

%%%%%%%%%%%%%%%%%%%%%%%%%%%%%%%%%%%%%%%%%%%%%%%%%%%%%%%%%%%%%%%%%%%%%%%%%%%
\section{The Measurement Chain}
%%%%%%%%%%%%%%%%%%%%%%%%%%%%%%%%%%%%%%%%%%%%%%%%%%%%%%%%%%%%%%%%%%%%%%%%%%%
\nsl{This chapter puts all building blocks of the lecture together. We
design a simple digital thermometer: a resistive temperature sensor
measures \SIrange{0}{100}{\celsius}, and the result is shown as a 3-bit
code on LEDs. Every stage of the chain is a circuit from an earlier
chapter: the bridge from chapter~4, the instrumentation amplifier from
chapter~5, the low-pass from chapter~6, the sample and hold from
chapter~7 and the flash converter from chapter~8. The LED with its series
resistor even goes back to chapter~1.
\newline
The design is done twice: on paper, with all calculations, and in
Falstad, where the complete chain runs and can be checked against the
paper values. The chapter is also the worked solution of last year's
capstone task (exercise~\ref{ex:capstone}).\newpage}

\begin{description}
\item[{\rot\bf Five stages from sensor to code}]~\\[-\lblskip]
\begin{center}
\resizebox{\textwidth}{!}{%
\begin{tikzpicture}[font=\sffamily\small,>=Stealth,x=1.15cm,
	b/.style={draw,thick,rounded corners,minimum height=15mm,text width=27mm,align=center,font=\sffamily\footnotesize}]
	\node[b,fill=rwucyan!10] (s1) at (0,0) {\textbf{1~Conditioning}\\bridge, INA,\\gain stage};
	\node[b,fill=rwucyan!20] (s2) at (3.4,0) {\textbf{2~Anti-aliasing}\\low-pass\\$f_c=\SI{2.34}{\hertz}$};
	\node[b,fill=rwuviolet!12] (s3) at (6.8,0) {\textbf{3 Sample \& Hold}\\$f_s=\SI{5}{\hertz}$};
	\node[b,fill=rwuviolet!22] (s4) at (10.2,0) {\textbf{4~Flash ADC}\\7~comparators,\\ladder};
	\node[b,fill=gray!12] (s5) at (13.6,0) {\textbf{5~Encoder \&}\\\textbf{Display}\\10~LEDs};
	\node[left=6mm of s1,align=center] (t) {Pt100\\\SIrange{0}{100}{\celsius}};
	\draw[->,thick] (t) -- (s1);
	\draw[->,thick] (s1) -- node[above,font=\sffamily\scriptsize]{0..4 V} (s2);
	\draw[->,thick] (s2) -- (s3);
	\draw[->,thick] (s3) -- node[above,font=\sffamily\scriptsize]{held} (s4);
	\draw[->,thick] (s4) -- node[above,font=\sffamily\scriptsize]{$C_1..C_7$} (s5);
	\node[below=3mm of s5,font=\sffamily\scriptsize] {$B_2B_1B_0$};
	\node[draw,thick,dashed,fill=yellow!15,align=center,font=\sffamily\scriptsize] (ref) at (6.8,-2.3)
		{reference \SI{4.00}{\volt}: bridge excitation and ladder (ratiometric)};
	\draw[->,dashed] (ref) -| (s1);
	\draw[->,dashed] (ref) -| (s4);
\end{tikzpicture}}
\end{center}
\bi
	\ite requirement: \SIrange{0}{100}{\celsius} $\rightarrow$ \SIrange{0}{4.00}{\volt} $\rightarrow$ code 000 \ldots 111
	\ite \SI{4}{\volt} full scale makes the step size round: $\LSB = \SI{4}{\volt}/2^3 = \SI{0.5}{\volt}$ $\widehat{=}$ \SI{12.5}{\celsius}
	\ite interference: \SI{50}{\hertz}, \SI{5}{\milli\volt} in series with the sensor (overhead lamps)
\ei
\os{\newpage}

\item[{\rot\bf Specification and design values}]~\\[-\lblskip]
\begin{center}
\begin{tabular}{lll}
\toprule
\textbf{Quantity} & \textbf{Requirement} & \textbf{Design value} \\
\midrule
sensor & resistive, slider in Falstad & Pt100, \SIrange{100}{138.5}{\ohm} \\
sensor current & $\leq \SI{1}{\milli\ampere}$ & \SIrange{1.000}{0.990}{\milli\ampere} \\
full scale & \SI{0}{\volt} \ldots \SI{4.00}{\volt} & $G = 10 \cdot 10.76 = 107.6$ \\
gain per stage & $< 100$ & 1, 10, 10.76 \\
anti-aliasing & $f_c \approx \SI{2.5}{\hertz}$ & \SI{6.8}{\kilo\ohm}, \SI{10}{\micro\farad}: \SI{2.34}{\hertz} \\
sampling & $f_s = \SI{5}{\hertz}$ & window \SI{20}{\milli\second}, $C_H = \SI{1}{\micro\farad}$ \\
ADC & 3 bit, LSB \SI{0.5}{\volt} & $8\times\SI{1}{\kilo\ohm}$, 7 comparators \\
display & LED + resistor & $10\times$ (LED + \SI{330}{\ohm}), \SI{9.8}{\milli\ampere} \\
\bottomrule
\end{tabular}
\end{center}
\falstad{capstone}
\os{\newpage}
\end{description}

%%%%%%%%%%%%%%%%%%%%%%%%%%%%%%%%%%%%%%%%%%%%%%%%%%%%%%%%%%%%%%%%%%%%%%%%%%%
\section{Stage 1: Signal Conditioning}
%%%%%%%%%%%%%%%%%%%%%%%%%%%%%%%%%%%%%%%%%%%%%%%%%%%%%%%%%%%%%%%%%%%%%%%%%%%
\nsl{The sensor is a Pt100 platinum resistor with
$R(T) = R_0\,(1+\alpha T)$, $R_0=\SI{100}{\ohm}$,
$\alpha = \SI{3.85e-3}{\per\kelvin}$, so it changes from \SI{100}{\ohm} at
\SI{0}{\celsius} to \SI{138.5}{\ohm} at \SI{100}{\celsius}. We feed the bridge
from the same \SI{4.00}{\volt} reference that later supplies the ADC ladder.
If the reference is \SI{1}{\percent} too high, the bridge output and all
thresholds of the converter are \SI{1}{\percent} too high as well, and the
code does not change: the measurement is \emph{ratiometric}. The reference
only needs to be stable while it is shared, not accurate.\newpage}

\begin{description}
\item[{\rot\bf Bridge design}]~\\[-\lblskip]
\begin{center}
\begin{circuitikz}[scale=0.85, transform shape]
  \draw (0,4) node[above]{$\Vref = \SI{4.00}{\volt}$} -- (3,4);
  \draw (0,4) to[R, l_=$R_1{=}\SI{3.9}{\kilo\ohm}$] (0,2) node[circ]{} node[left]{$V_S$}
        to[thR, l_=Pt100] (0,0.5)
        to[sV, l_=\SI{5}{\milli\volt}\,\SI{50}{\hertz}] (0,-1) node[ground]{};
  \draw (3,4) to[R, l=$R_3{=}\SI{3.9}{\kilo\ohm}$] (3,2) node[circ]{} node[right]{$V_R$}
        to[R, l=$R_4{=}\SI{100}{\ohm}$] (3,-1) node[ground]{};
  \draw (0,2) to[open, v^>=$V_{BR}$] (3,2);
\end{circuitikz}
\end{center}
\bi
	\ite sensor current $I_S = \Vref/(R_1+R) = \SI{4}{\volt}/\SI{4000}{\ohm} = \SI{1.000}{\milli\ampere}$ at \SI{0}{\celsius} (\SI{0.990}{\milli\ampere} at \SI{100}{\celsius})
	\ite balanced at \SI{0}{\celsius}: $R_4 = R(\SI{0}{\celsius})$, so $V_{BR}=0$
	\ite $V_{BR}(T) = \Vref\left(\frac{R(T)}{R_1+R(T)} - \frac{R_4}{R_3+R_4}\right)$,
	\quad $V_{BR}(\SI{100}{\celsius}) = \SI{37.18}{\milli\volt}$
\ei
\nsl{The bridge output is not exactly linear in $T$, because the sensor
current falls slightly with rising $R$. With $R_1 = 39\,R_0$ the error is
small: at \SI{50}{\celsius} the output is \SI{0.09}{\milli\volt} higher than
the straight line, \SI{0.25}{\percent} of full scale, far below one LSB
(\SI{12.5}{\percent}). A larger $R_1$ would reduce it further, but also the
signal.}
\os{\newpage}

\item[{\rot\bf Instrumentation amplifier and gain stage}]~\\[-\lblskip]
\begin{center}
\begin{circuitikz}[scale=0.72, transform shape]
  % buffers: V_R on top (to the inverting side), V_S below
  \node[op amp, noinv input up] (b1) at (1.2,3) {};
  \draw (b1.+) to[short,-o] ++(-0.6,0) node[left]{$V_R$};
  \draw (b1.-) -- ++(0,-0.6) -| (b1.out);
  \node[op amp, noinv input up] (b2) at (1.2,-0.6) {};
  \draw (b2.+) to[short,-o] ++(-0.6,0) node[left]{$V_S$};
  \draw (b2.-) -- ++(0,-0.6) -| (b2.out);
  % difference amplifier, G1 = 10
  \node[op amp] (da) at (7,1.2) {};
  \draw (b1.out) -- ++(0.4,0) |- ($(da.-)+(-2.4,0)$) to[R, l=\SI{10}{\kilo\ohm}] (da.-);
  \draw (b2.out) -- ++(0.4,0) |- ($(da.+)+(-2.4,0)$) to[R, l_=\SI{10}{\kilo\ohm}] (da.+);
  \draw ($(da.-)+(-0.4,0)$) node[circ]{} -- ++(0,1.2) coordinate(t)
        to[R, l=\SI{100}{\kilo\ohm}] (t -| da.out) -- (da.out);
  \draw ($(da.+)+(-0.4,0)$) node[circ]{} to[R, l=\SI{100}{\kilo\ohm}] ++(0,-1.8) node[ground]{};
  % gain stage, G2 = 10.76
  \node[op amp, noinv input up, anchor=+] (gs) at ($(da.out)+(1.8,0)$) {};
  \draw (da.out) -- (gs.+) node[pos=0.4,above]{\scriptsize $V_\mathrm{DIFF}$};
  \draw (gs.-) -- ++(0,-0.9) coordinate(m) to[R, l_=\SI{10}{\kilo\ohm}] ++(0,-1.6) node[ground]{};
  \draw (m) node[circ]{} to[R, l_=\SI{97.6}{\kilo\ohm}] ++(2.6,0) coordinate(o2) -- (o2 |- gs.out);
  \draw (gs.out) -- (o2 |- gs.out) node[circ]{} to[short,-o] ++(0.7,0) node[right]{$V_\mathrm{COND}$};
\end{circuitikz}
\end{center}
\bi
	\ite buffers: input current zero, the bridge is not loaded
	\ite difference amplifier: $V_\mathrm{DIFF} = \frac{\SI{100}{\kilo\ohm}}{\SI{10}{\kilo\ohm}}(V_S - V_R) = 10\,V_{BR}$
	\ite gain stage: $G_2 = 1 + \frac{\SI{97.6}{\kilo\ohm}}{\SI{10}{\kilo\ohm}} = 10.76$
\ei
\os{\newpage}

\item[{\rot\bf Gain: how much, and in how many stages?}]~\\[-\lblskip]
\bi
	\ite required total gain $G = \SI{4.00}{\volt}/\SI{37.18}{\milli\volt} = 107.6$
	\ite split: $G_1 = 10$ (difference amplifier), $G_2 = 4.00/(10 \cdot 0.03718) = 10.758$
	\itee E96: $R_F = \SI{97.6}{\kilo\ohm}$ $\Rightarrow$ $G_2 = 10.76$, full scale \SI{4.0005}{\volt}
	\itee in hardware: $R_F = \SI{91}{\kilo\ohm} + \SI{10}{\kilo\ohm}$ trimmer, adjusted at \SI{100}{\celsius}
	\ite finite open-loop gain $A = 10^5$: relative gain error $\approx G/A$
	\itee one stage $G = 107.6$: \SI{0.11}{\percent}; two stages: $2 \times \SI{0.01}{\percent}$
	\itee a single stage with $G > 1000$ would lose more than \SI{1}{\percent}: keep $G < 100$ per stage
\ei
\keyeq{Conditioning}{$\displaystyle V_\mathrm{COND} = G_1 G_2 \, \Vref\left(\frac{R(T)}{R_1+R(T)}-\frac{R_4}{R_3+R_4}\right)$,\quad $G_1G_2 = 107.6$}
\nsl{Splitting the gain over two stages has a second advantage: the zero
point is set by the bridge ($R_4$) and the span by one resistor ($R_F$).
The two adjustments do not influence each other, so the trimming procedure
is simple: set the sensor to \SI{100}{\ohm} and check \SI{0}{\volt}, then set
it to \SI{138.5}{\ohm} and adjust $R_F$ for \SI{4.00}{\volt}.}
\os{\newpage}
\end{description}

%%%%%%%%%%%%%%%%%%%%%%%%%%%%%%%%%%%%%%%%%%%%%%%%%%%%%%%%%%%%%%%%%%%%%%%%%%%
\section{Stage 2: Anti-Aliasing}
%%%%%%%%%%%%%%%%%%%%%%%%%%%%%%%%%%%%%%%%%%%%%%%%%%%%%%%%%%%%%%%%%%%%%%%%%%%
\nsl{The temperature changes slowly, so five samples per second are
plenty. The sampling theorem then limits the usable bandwidth to
$f_s/2 = \SI{2.5}{\hertz}$. Everything above this frequency must be removed
\emph{before} the sampler, because after sampling it can no longer be
distinguished from a real signal. The \SI{50}{\hertz} interference is the
dangerous part: after the gain of 107.6 it has an amplitude of about half
a volt, one full LSB.\newpage}

\begin{description}
\item[{\rot\bf Where does the 50\,Hz go?}]~\\[-\lblskip]
\bi
	\ite sampling with $f_s$ folds a frequency $f$ to $f_\mathrm{alias} = |f - k f_s|$ (nearest $k$)
	\ite $f = \SI{50}{\hertz}$, $f_s = \SI{5}{\hertz}$: $|50 - 10\cdot 5| = \SI{0}{\hertz}$
	\itee every sample sees the interference at the same phase
	\itee the result is a \textbf{constant offset}: nothing on the scope looks wrong
	\ite amplitude at $V_\mathrm{COND}$: $\SI{5}{\milli\volt}\cdot\frac{R_1}{R_1+R}\cdot 107.6 = \SI{0.52}{\volt}$ (about 1 LSB)
\ei
\begin{center}
\begin{tikzpicture}
\begin{axis}[width=0.8\textwidth,height=38mm,xmin=0,xmax=0.42,ymin=-1.3,ymax=1.3,
	xlabel={$t$ / s},ytick={-1,0,1},xtick={0,0.1,0.2,0.3,0.4},
	xticklabel style={/pgf/number format/fixed},axis lines=left,font=\scriptsize]
	\addplot[rwucyan,samples=800,domain=0:0.42] {cos(deg(2*pi*50*x))};
	\addplot[rwuviolet,only marks,mark=*,mark size=1.6pt] coordinates {(0.02,1) (0.22,1) (0.42,1)};
	\addplot[rwuviolet,thick,dashed,domain=0:0.42] {1};
\end{axis}
\end{tikzpicture}
\end{center}
\os{\newpage}

\item[{\rot\bf Anti-aliasing low-pass}]~\\[-\lblskip]
\begin{center}
\begin{circuitikz}[scale=0.85, transform shape]
  \draw (0,0) node[left]{$V_\mathrm{COND}$} to[short,o-] (0.3,0) to[R, l=\SI{6.8}{\kilo\ohm}] (2.5,0) coordinate(f)
        to[C, l=\SI{10}{\micro\farad}] (2.5,-1.8) node[ground]{};
  \draw (4,-0.48) node[op amp, noinv input up, anchor=+] (lp) {};
  \draw (f) -- (lp.+);
  \draw (lp.-) |- ++(0.5,-0.9) -| (lp.out);
  \draw (lp.out) to[short,-o] ++(0.6,0) node[right]{$V_\mathrm{AA}$};
\end{circuitikz}
\end{center}
\bi
	\ite $f_c = \frac{1}{2\pi RC} = \frac{1}{2\pi\cdot\SI{6.8}{\kilo\ohm}\cdot\SI{10}{\micro\farad}} = \SI{2.34}{\hertz} \leq f_s/2$
	\ite at \SI{50}{\hertz}: $|H| = 1/\sqrt{1+(50/2.34)^2} = 0.047$ ($-\SI{26.6}{\dB}$)
	\ite residual offset: $\SI{0.52}{\volt}\cdot 0.047 = \SI{25}{\milli\volt} < \LSB/2 = \SI{250}{\milli\volt}$
	\ite the follower decouples the RC network from the S\&H input
\ei
\nsl{For 3 bits a first-order filter is sufficient. For a 10-bit converter
($\LSB/2 = \SI{2}{\milli\volt}$) the residual \SI{25}{\milli\volt} would be
twelve times too large; a second- or third-order filter (chapter~6), a higher
sampling rate followed by digital averaging, or a sampling rate that is
\emph{not} a sub-multiple of \SI{50}{\hertz} would then be needed.}
\designcard{Nyquist \& anti-aliasing (capstone values)}{%
$f_s = \SI{5}{\hertz}$ $\Rightarrow$ bandwidth $\leq \SI{2.5}{\hertz}$;
\quad $f_\mathrm{alias} = |f - k f_s|$; \quad
$f_c = 1/(2\pi RC) \leq f_s/2$; \quad residual $= \hat v\,/\sqrt{1+(f/f_c)^2}$ must stay $< \LSB/2$.}
\os{\newpage}
\end{description}

%%%%%%%%%%%%%%%%%%%%%%%%%%%%%%%%%%%%%%%%%%%%%%%%%%%%%%%%%%%%%%%%%%%%%%%%%%%
\section{Stage 3: Sample and Hold}
%%%%%%%%%%%%%%%%%%%%%%%%%%%%%%%%%%%%%%%%%%%%%%%%%%%%%%%%%%%%%%%%%%%%%%%%%%%
\nsl{The flash converter is fast and would not strictly need a sample and
hold. We use one anyway, because it defines the sampling instant: the
input of the converter changes only once per sampling period, so the LEDs
show a stable code, and the aliasing experiment becomes visible. The
circuit is the buffered S\&H of chapter~7.\newpage}

\begin{description}
\item[{\rot\bf Buffered sample and hold}]~\\[-\lblskip]
\begin{center}
\begin{circuitikz}[scale=0.8, transform shape]
  \draw (0,-0.48) node[op amp, noinv input up, anchor=+] (i) {};
  \draw (i.-) |- ++(0.5,-0.9) -| (i.out);
  \draw (i.+) to[short,-o] ++(-0.5,0) node[left]{$V_\mathrm{AA}$};
  \draw (i.out) to[nos, l={CLK}, a={FET switch}] ++(2.2,0) coordinate(h)
        to[C, l=$C_H$] ++(0,-1.8) node[ground]{};
  \draw (h) -- ++(1.2,0) coordinate(hh);
  \draw (hh) ++(0.0,0) node[op amp, noinv input up, anchor=+] (o) {};
  \draw (o.-) |- ++(0.5,-0.9) -| (o.out);
  \draw (o.out) to[short,-o] ++(0.5,0) node[right]{$V_\mathrm{HOLD}$};
\end{circuitikz}
\end{center}
\bi
	\ite clock \SI{5}{\hertz}, duty \SI{10}{\percent}: switch closed for \SI{20}{\milli\second} (track), open for \SI{180}{\milli\second} (hold)
	\ite input buffer: charges $C_H$ quickly; output buffer: no load on $C_H$
\ei
\os{\newpage}

\item[{\rot\bf Choosing the hold capacitor}]~\\[-\lblskip]
\bi
	\ite \textbf{charging:} $\tau = R_\mathrm{on} C_H = \SI{20}{\ohm}\cdot\SI{1}{\micro\farad} = \SI{20}{\micro\second}$
	\itee settle to $\LSB/2$ of 3 bits: $t = \tau\ln 2^{4} = 2.8\,\tau = \SI{55}{\micro\second} \ll \SI{20}{\milli\second}$
	\ite \textbf{droop:} leakage $I_L = \SI{1}{\nano\ampere}$ (switch + buffer input)
	\itee $\Delta V = I_L\, t_\mathrm{hold}/C_H = \SI{1}{\nano\ampere}\cdot\SI{0.18}{\second}/\SI{1}{\micro\farad} = \SI{0.18}{\milli\volt}$
	\ite both limits leave a wide range: \SI{10}{\nano\farad} \ldots \SI{100}{\micro\farad} would work
	\itee choose a film capacitor with low leakage and dielectric absorption: $C_H = \SI{1}{\micro\farad}$
\ei
\nsl{In hardware a small resistor (e.g.\ \SI{1}{\kilo\ohm}, as in the
examples of chapter~7) is usually put in series with the switch: it limits
the charging current that the input buffer must deliver and isolates the
buffer from the capacitive load. Then $\tau = \SI{1.02}{\milli\second}$
and settling to $\LSB/2$ takes \SI{2.8}{\milli\second}, still well inside
the window. The limits are far apart because the signal is slow and the
resolution low. They move together quickly in a fast, high-resolution
system: a 12-bit converter sampling at \SI{1}{\mega\hertz} has a track time
of a few hundred nanoseconds and needs settling to $2^{-13}$, which is why
real S\&H stages use picofarad capacitors and low-resistance switches.}
\os{\newpage}
\end{description}

%%%%%%%%%%%%%%%%%%%%%%%%%%%%%%%%%%%%%%%%%%%%%%%%%%%%%%%%%%%%%%%%%%%%%%%%%%%
\section{Stages 4 and 5: Flash ADC, Encoder and Display}
%%%%%%%%%%%%%%%%%%%%%%%%%%%%%%%%%%%%%%%%%%%%%%%%%%%%%%%%%%%%%%%%%%%%%%%%%%%
\nsl{The held voltage is compared with seven thresholds at once. The
thresholds come from a ladder of eight equal resistors across the
\SI{4.00}{\volt} reference. Each comparator answers one question, ``is
the input above $k \cdot \SI{0.5}{\volt}$?'', and together they form a
thermometer code: all comparators below the input voltage output 1, all
above output 0. An encoder turns the seven bits into the 3-bit binary
number.\newpage}

\begin{description}
\item[{\rot\bf Reference ladder and comparators}]~\\[-\lblskip]
\begin{center}
\begin{circuitikz}[scale=0.62, transform shape]
  \draw (0,6) node[above]{\SI{4.00}{\volt}} to[R, l_=\SI{1}{\kilo\ohm}] (0,4.5) coordinate(t7)
        to[R, l_=\SI{1}{\kilo\ohm}] (0,3) coordinate(t6);
  \draw[dashed] (0,3) -- (0,1.5);
  \draw (0,1.5) coordinate(t1x) to[R, l_=\SI{1}{\kilo\ohm}] (0,0) coordinate(t1)
        to[R, l_=\SI{1}{\kilo\ohm}] (0,-1.5) node[ground]{};
  \foreach \t/\k/\v in {t7/7/3.5, t6/6/3.0, t1/1/0.5} {
    \draw (\t) node[circ]{} -- ++(2.2,0) coordinate(m\k);
    \draw (m\k) ++(1.2,0.48) node[op amp, noinv input up] (c\k) {};
    \draw (m\k) -- (c\k.-);
    \draw (c\k.+) -- ++(-0.4,0) node[left]{\scriptsize $V_\mathrm{HOLD}$};
    \draw (c\k.out) to[short,-o] ++(0.5,0) node[right]{$C_\k$};
    \node[left] at (\t) {\scriptsize \SI{\v}{\volt}};
  }
\end{circuitikz}
\end{center}
\bi
	\ite ladder current $\SI{4}{\volt}/\SI{8}{\kilo\ohm} = \SI{0.5}{\milli\ampere}$; comparator inputs draw no current
	\ite $C_k = 1$ if $V_\mathrm{HOLD} > k\cdot\SI{0.5}{\volt}$; the number of ones is the level 0\ldots7
\ei
\os{\newpage}

\item[{\rot\bf Quantization and encoder}]~\\[-\lblskip]
\begin{center}
\small
\begin{tabular}{cc|ccccccc|ccc}
\toprule
$V_\mathrm{HOLD}$ / V & level & $C_7$ & $C_6$ & $C_5$ & $C_4$ & $C_3$ & $C_2$ & $C_1$ & $B_2$ & $B_1$ & $B_0$ \\
\midrule
0.0--0.5 & 0 & 0&0&0&0&0&0&0 & 0&0&0 \\
0.5--1.0 & 1 & 0&0&0&0&0&0&1 & 0&0&1 \\
1.0--1.5 & 2 & 0&0&0&0&0&1&1 & 0&1&0 \\
1.5--2.0 & 3 & 0&0&0&0&1&1&1 & 0&1&1 \\
2.0--2.5 & 4 & 0&0&0&1&1&1&1 & 1&0&0 \\
2.5--3.0 & 5 & 0&0&1&1&1&1&1 & 1&0&1 \\
3.0--3.5 & 6 & 0&1&1&1&1&1&1 & 1&1&0 \\
3.5--4.0 & 7 & 1&1&1&1&1&1&1 & 1&1&1 \\
\bottomrule
\end{tabular}
\end{center}
\bi
	\ite $B_2 = C_4$ \quad (level $\geq 4$)
	\ite $B_1 = C_2\overline{C_4} + C_6$ \quad (levels 2, 3, 6, 7)
	\ite $B_0 = C_1\overline{C_2} + C_3\overline{C_4} + C_5\overline{C_6} + C_7$ \quad (odd levels: the edge of the thermometer is at an odd position)
\ei
\os{\newpage}

\item[{\rot\bf Display}]~\\[-\lblskip]
\bi
	\ite 7 red LEDs from $C_1\ldots C_7$: bar graph, a thermometer you can read without decoding
	\ite 3 green LEDs from $B_2B_1B_0$: the binary result
	\ite each LED with its own series resistor (chapter~1), driven by a \SI{5}{\volt} logic output:
	$R = \frac{\SI{5}{\volt}-V_F}{I_F} = \frac{\SI{5}{\volt}-\SI{1.8}{\volt}}{\SI{10}{\milli\ampere}} = \SI{320}{\ohm}$ $\rightarrow$ E12: \SI{330}{\ohm}
	\itee Falstad: $I_F = \SI{9.8}{\milli\ampere}$, $V_F = \SI{1.78}{\volt}$
	\ite never one common resistor for several LEDs: the current would depend on how many are on
\ei
\os{\newpage}
\end{description}

%%%%%%%%%%%%%%%%%%%%%%%%%%%%%%%%%%%%%%%%%%%%%%%%%%%%%%%%%%%%%%%%%%%%%%%%%%%
\section{The Complete Chain in Falstad}
%%%%%%%%%%%%%%%%%%%%%%%%%%%%%%%%%%%%%%%%%%%%%%%%%%%%%%%%%%%%%%%%%%%%%%%%%%%
\nsl{The Falstad model contains all five stages in labelled boxes. The
stages are connected with labeled nodes (\texttt{VCOND}, \texttt{VAA},
\texttt{VHOLD}, \texttt{C1}\ldots\texttt{C7}, \texttt{B0}\ldots\texttt{B2});
nodes with the same name are connected, which keeps the drawing readable.
The sensor resistance is set with the slider ``Sensor R''. The first scope
shows the filtered signal and the held staircase, the second the
conditioned signal and the held value.\newpage}

\begin{description}
\item[{\rot\bf Using the model}]~\\[-\lblskip]
\bi
	\ite open \falstad{capstone}
	\ite move the slider from \SI{100}{\ohm} to \SI{138.5}{\ohm}: the bar graph fills up, the binary LEDs count
	\ite hover over a label to read its voltage; the held value changes every \SI{200}{\milli\second}
	\ite after a slider step, wait about \SI{0.5}{\second}: the low-pass needs $5\tau = 5\cdot\SI{68}{\milli\second}$
\ei
\os{\newpage}

\item[{\rot\bf Verification table}]~\\[-\lblskip]
\nsl{For each test point the sensor resistance is computed from the
conditioning equation so that $V_\mathrm{COND}$ lies in the centre of a code
bin. The ``Falstad'' columns are the values of the simulation (with the
\SI{50}{\hertz} interference switched on) just before the next sample.}
\begin{center}
\begin{tabular}{rrrrcc}
\toprule
$T$ / \si{\celsius} & $R$ / \si{\ohm} & $V_\mathrm{COND}$ target & $V_\mathrm{HOLD}$ Falstad & code expected & code Falstad \\
\midrule
 6.2 & 102.38 & \SI{0.25}{\volt} & \SI{0.251}{\volt} & 000 & 000 \\
18.6 & 107.16 & \SI{0.75}{\volt} & \SI{0.751}{\volt} & 001 & 001 \\
43.5 & 116.75 & \SI{1.75}{\volt} & \SI{1.749}{\volt} & 011 & 011 \\
56.0 & 121.56 & \SI{2.25}{\volt} & \SI{2.250}{\volt} & 100 & 100 \\
81.1 & 131.22 & \SI{3.25}{\volt} & \SI{3.249}{\volt} & 110 & 110 \\
93.7 & 136.07 & \SI{3.75}{\volt} & \SI{3.750}{\volt} & 111 & 111 \\
\bottomrule
\end{tabular}
\end{center}
\bi
	\ite the temperatures differ slightly from the linear mapping \SI{25}{\celsius\per\volt} (6, 19, 44, 56, 81, \SI{94}{\celsius}) because of the bridge nonlinearity
\ei
\os{\newpage}

\item[{\rot\bf The aliasing experiment}]~\\[-\lblskip]
\bi
	\ite \falstad{capstone_nofilter}: the S\&H takes $V_\mathrm{COND}$ directly
	\ite $V_\mathrm{COND}$ carries \SI{0.52}{\volt} of \SI{50}{\hertz}; the S\&H always samples at the same phase
\ei
\begin{center}
\begin{tabular}{lccc}
\toprule
 & $V_\mathrm{COND}$ (DC) & $V_\mathrm{HOLD}$ & code \\
\midrule
without filter & \SI{1.75}{\volt} & \SI{2.27}{\volt} & 100 (wrong) \\
without filter & \SI{0.25}{\volt} & \SI{0.77}{\volt} & 001 (wrong) \\
with filter & \SI{1.75}{\volt} & \SI{1.75}{\volt} & 011 \\
\bottomrule
\end{tabular}
\end{center}
\bi
	\ite without the filter the held value is a clean staircase, just \SI{0.52}{\volt} too high: one full code too high
	\ite the size of the offset depends on the phase between mains and sampling clock (here the worst case)
\ei
\os{\newpage}
\end{description}

%%%%%%%%%%%%%%%%%%%%%%%%%%%%%%%%%%%%%%%%%%%%%%%%%%%%%%%%%%%%%%%%%%%%%%%%%%%
\section{Exercises}
%%%%%%%%%%%%%%%%%%%%%%%%%%%%%%%%%%%%%%%%%%%%%%%%%%%%%%%%%%%%%%%%%%%%%%%%%%%

\exercise{From sensor to digital value (capstone task 2025/26)}\label{ex:capstone}
Build a complete data-acquisition chain in Falstad that turns a slowly
changing sensor signal (\SIrange{0}{100}{\celsius}) into a 3-bit code. Label
the five stages Conditioning~| Anti-aliasing~| Sample \& Hold~| Flash ADC~|
Encoder \& Display.
\begin{talist}
\ta Conditioning: resistive sensor in a Wheatstone bridge, sensor current
	$\leq\SI{1}{\milli\ampere}$, sensor as a resistor with a slider;
	instrumentation-amplifier front end (two unity-gain buffers and a
	difference amplifier); gain trimmed to \SI{0}{\volt} at the bottom and
	exactly \SI{4.00}{\volt} at the top of the range; less than 100 per stage.
\ta Anti-aliasing: sampling clock $f_s=\SI{5}{\hertz}$. Add the
	\SI{50}{\hertz}, \SI{5}{\milli\volt} noise source in series with the
	sensor. Observe the held value without a filter, then add an active
	low-pass with $f_c\approx\SI{2.5}{\hertz}$ and justify $f_c$ from the
	Nyquist criterion.
\ta Sample and hold at \SI{5}{\hertz}: follower, FET switch, hold capacitor,
	follower. Choose the hold capacitor.
\ta Flash ADC: ladder $8\times\SI{1}{\kilo\ohm}$ from \SI{4}{\volt}, seven
	comparators, thermometer code.
\ta Encoder ($B_2 = C_4$, $B_1 = C_2\overline{C_4} + C_6$,
	$B_0 = C_1\overline{C_2} + C_3\overline{C_4} + C_5\overline{C_6} + C_7$)
	and display: 7-LED bar graph and three binary LEDs, each LED with its own
	resistor.
\ta Verification table for six inputs near the centres of different code
	bins, and one line on the aliasing experiment.
\end{talist}

\begin{solution}
The design of this chapter is one possible solution; the complete model
is \falstad{capstone}.
\begin{talist}
\ta Pt100 with $R_1=R_3=\SI{3.9}{\kilo\ohm}$, $R_4=\SI{100}{\ohm}$ at
	$\Vref=\SI{4.00}{\volt}$: $I_S = \SI{1.000}{\milli\ampere}$, bridge
	output \SIrange{0}{37.18}{\milli\volt}. Buffers, difference amplifier
	$G_1=10$, non-inverting stage $G_2 = 1+\SI{97.6}{\kilo\ohm}/\SI{10}{\kilo\ohm}=10.76$:
	$V_\mathrm{COND} = \SIrange{0}{4.00}{\volt}$.
\ta Without filter the 50\,Hz aliases to
	$|50-10\cdot5| = \SI{0}{\hertz}$; the held value is offset by up to
	\SI{0.52}{\volt} (Falstad: \SI{1.75}{\volt} becomes \SI{2.27}{\volt},
	code 100 instead of 011, \falstad{capstone_nofilter}). Filter
	$R=\SI{6.8}{\kilo\ohm}$, $C=\SI{10}{\micro\farad}$, $f_c=\SI{2.34}{\hertz}\leq f_s/2$,
	50\,Hz attenuated to \SI{4.7}{\percent}: \SI{25}{\milli\volt} residual,
	below $\LSB/2$.
\ta $C_H=\SI{1}{\micro\farad}$: $\tau = \SI{20}{\micro\second}$ charges
	within the \SI{20}{\milli\second} window; droop
	\SI{0.18}{\milli\volt} per hold period for \SI{1}{\nano\ampere} leakage.
\ta Taps at $0.5, 1.0, \ldots, \SI{3.5}{\volt}$; ladder current
	\SI{0.5}{\milli\ampere}.
\ta 3 inverters, 4 AND gates, 2 OR gates (one with 4 inputs);
	10 LEDs with \SI{330}{\ohm}, \SI{9.8}{\milli\ampere} each.
\ta See the verification table above. Aliasing line: ``Without the
	anti-aliasing filter the held value showed a constant offset of
	\SI{+0.52}{\volt} (one code too high); with the \SI{2.34}{\hertz}
	low-pass the offset fell below \SI{30}{\milli\volt} and all codes were
	correct.''
\end{talist}
\end{solution}

\exercise{Ratiometric measurement}
The \SI{4.00}{\volt} reference of the capstone is in fact \SI{4.10}{\volt}.
\begin{talist}
\ta How large is $V_\mathrm{COND}$ at \SI{100}{\celsius}?
\ta Which code does the converter show at \SI{100}{\celsius} and at
	\SI{56}{\celsius}? Compare with a design in which the bridge is supplied
	from a separate, exact \SI{4.00}{\volt} source.
\end{talist}

\begin{solution}
\begin{talist}
\ta The bridge output scales with the reference:
	$V_\mathrm{COND} = \SI{4.0005}{\volt}\cdot 4.10/4.00 = \SI{4.10}{\volt}$.
\ta The thresholds scale as well ($k\cdot\SI{0.5125}{\volt}$). The ratio
	$V_\mathrm{COND}/\Vref$ is unchanged, so both codes are unchanged: 111 at
	\SI{100}{\celsius}, 100 at \SI{56}{\celsius} ($V_\mathrm{COND} =
	\SI{2.25}{\volt}\cdot1.025 = \SI{2.306}{\volt}$, threshold of level 4:
	\SI{2.05}{\volt}, of level 5: \SI{2.5625}{\volt}). With an exact bridge
	supply and a \SI{4.10}{\volt} ladder the thresholds would be
	\SI{2.5}{\percent} too high while the signal is correct; a value just
	above a threshold (e.g.\ \SI{2.02}{\volt}) would then be read one code
	too low. Ratiometric operation removes the reference from the error
	budget.
\end{talist}
\end{solution}

\exercise{Another sampling rate}
The sampling clock of the capstone is changed to $f_s = \SI{4}{\hertz}$
and the anti-aliasing filter is removed.
\begin{talist}
\ta To which frequency does the \SI{50}{\hertz} interference alias?
\ta Describe what the LEDs show for a constant temperature with
	$V_\mathrm{COND}=\SI{1.75}{\volt}$.
\ta Is this better or worse than the \SI{0}{\hertz} alias at \SI{5}{\hertz}?
\end{talist}

\begin{solution}
\begin{talist}
\ta $|50 - 12\cdot 4| = \SI{2}{\hertz}$, exactly the Nyquist frequency
	$f_s/2$: consecutive samples see the interference with a phase step of
	$2\pi\cdot 50/4 = 25\pi$, so its sign alternates from sample to sample.
\ta The held value alternates between $1.75 + \hat v\sin\varphi$ and
	$1.75 - \hat v\sin\varphi$ ($\hat v = \SI{0.52}{\volt}$): the LEDs
	blink at \SI{2}{\hertz} between two codes, e.g.\ 100 and 010 in the
	worst case, or stay at 011 if the phase happens to be $\varphi \approx 0$.
\ta It is visible, so it is ``better'' than a silent offset, but the
	reading is still wrong. Only the low-pass in front of the sampler
	removes the error for every sampling rate.
\end{talist}
\end{solution}

\exercise{Four bits}
The capstone shall be extended to 4 bits with the same \SI{4.00}{\volt}
full scale.
\begin{talist}
\ta What is the new LSB, in volts and in \si{\celsius}? How many
	comparators and ladder resistors are needed?
\ta Is the first-order anti-aliasing filter still sufficient?
\ta Is the bridge nonlinearity still acceptable?
\end{talist}

\begin{solution}
\begin{talist}
\ta $\LSB = \SI{4}{\volt}/16 = \SI{0.25}{\volt}$ $\widehat{=}$
	\SI{6.25}{\celsius}; $2^4-1 = 15$ comparators, 16 resistors (taps at
	$k\cdot\SI{0.25}{\volt}$).
\ta Residual \SI{25}{\milli\volt} $< \LSB/2 = \SI{125}{\milli\volt}$: yes,
	still a factor five of margin.
\ta The maximum deviation from the straight line is about
	\SI{0.09}{\milli\volt}$\cdot 107.6 = \SI{10}{\milli\volt}$
	$\ll \SI{125}{\milli\volt}$: yes.
\end{talist}
\end{solution}

\exercise{Hold capacitor at the limits}
In the capstone S\&H the analog switch has $R_\mathrm{on}=\SI{1}{\kilo\ohm}$
and the total leakage current is \SI{10}{\nano\ampere}. Determine the
smallest and the largest usable hold capacitor if the droop during one
hold period and the remaining charging error at the end of the window
shall both stay below $\LSB/4 = \SI{125}{\milli\volt}$ for a full-scale
step.

\begin{solution}
Droop: $I_L\,t_\mathrm{hold}/C_H < \SI{0.125}{\volt}$ $\Rightarrow$
$C_H > \SI{10}{\nano\ampere}\cdot\SI{0.18}{\second}/\SI{0.125}{\volt} = \SI{14.4}{\nano\farad}$.
Charging: $\SI{4}{\volt}\cdot e^{-t_w/\tau} < \SI{0.125}{\volt}$
$\Rightarrow$ $t_w > \tau\ln 32 = 3.47\,\tau$, so
$\tau < \SI{20}{\milli\second}/3.47 = \SI{5.8}{\milli\second}$ and
$C_H < \SI{5.8}{\milli\second}/\SI{1}{\kilo\ohm} = \SI{5.8}{\micro\farad}$.
Any value between about \SI{15}{\nano\farad} and \SI{5.8}{\micro\farad}
works; \SI{1}{\micro\farad} is a good choice in the middle (on a
logarithmic scale). In Falstad set the switch's on resistance to
\SI{1}{\kilo\ohm} and $C_H$ to \SI{10}{\micro\farad} in
\falstad{capstone}: after a slider step the held staircase needs several
samples to reach the new value.
\end{solution}

%%% Local Variables:
%%% mode: latex
%%% TeX-master: "docu"
%%% End:
